\chapter{ระเบียบวิธีวิจัย}

\section{ภาพรวมของแบบจำลอง}

แบบจำลองที่นำเสนอใช้ Transformer เพื่อเข้ารหัสหน้าต่างข้อมูล OHLCV หลายหลักทรัพย์ จากนั้นสร้างข้อมูลกลับคืนและประมาณพารามิเตอร์ของการแจกแจงผลตอบแทนแบบ GBM anomaly score ประกอบด้วยองค์ประกอบจาก reconstruction error, negative log-likelihood และ divergence ระหว่างการแจกแจงผลตอบแทนที่สังเกตได้กับการแจกแจงที่แบบจำลองคาดหมาย

\section{Problem Formulation}

ให้ $X_{t:t+w}$ แทนหน้าต่างข้อมูล OHLCV ที่มีความยาว $w$ สำหรับหลายหลักทรัพย์ เป้าหมายคือเรียนรู้ฟังก์ชันให้คะแนน $s(X_{t:t+w})$ ซึ่งมีค่าสูงเมื่อหน้าต่างดังกล่าวมีพฤติกรรมเบี่ยงเบนจากรูปแบบปกติ

\section{GBM-Implied Return Consistency}

ภายใต้กรอบ GBM ผลตอบแทน log-return สามารถเชื่อมโยงกับ drift และ volatility ได้ งานวิจัยนี้ใช้แนวคิดดังกล่าวเพื่อสร้างตัวชี้วัดความไม่สอดคล้องระหว่างผลตอบแทนที่สังเกตได้และการแจกแจงที่แบบจำลองคาดหมาย

\section{Dynamic Thresholding}

นอกจาก anomaly score โดยตรง งานวิจัยนี้พิจารณา score-change dynamic threshold เพื่อจับการเปลี่ยนแปลงฉับพลันของคะแนนความผิดปกติ กฎนี้เหมาะสำหรับการแจ้งเตือนที่ต้องการ precision สูง แต่ยังต้องวิเคราะห์ sensitivity และ lead-time เพิ่มเติมก่อนใช้เป็นข้อสรุปเชิง early warning

